\section{Conclusion}

Staple dispersion supplies a gauge-invariant answer to a concrete local
question: how much incident geodesic-quadratic energy survives after the
central link is optimized?  On a coherent nonlinear chart it is equivalent,
up to the spherical factor $(\sin r/r)^2$, to pairwise relational dispersion.
In the linearized finite block its probability law retains the exact
shared-link dependence and reduces to five certified Gaussian modes.

The reverse problem is more revealing.  Prescribing those five relative modes
does not define a local operation on links.  Every nonzero lift creates
external curvature, with a certified cost of at least $37/33$ times the
desired squared displacement on the test block.  The result is modest in
scale but sharp in meaning: what looks like a local contraction after a
representation change can be incompatible with microscopic locality.  Any
larger constructive program using such a repair must control that
incompatibility rather than treating the staple coordinates as independent.
